Una masa de $100\,\mathrm{kg}$ está unida en paralelo a un muelle de $80\,\mathrm{N/m}$ y a un amortiguador de $40\,\mathrm{N\,s/m}$. Se mide $x$ en metros, positivo hacia abajo, desde el equilibrio previo cargado por la gravedad. En $t=0$ se aplica una fuerza adicional constante de $15\,\mathrm N$, partiendo del reposo y con $x(0)=0$. Use segundos. Para estudiar la gráfica, se proporciona la solución ya obtenida del PVI $100x''+40x'+80x=15$: $x(t)=\frac3{16}\left[1-e^{-t/5}\left(\cos\frac{\sqrt{19}t}{5}+\frac1{\sqrt{19}}\sin\frac{\sqrt{19}t}{5}\right)\right]$, para $t\ge0$.

\begin{enumerate}[label=(\alph*),ref=\alph*,leftmargin=*]
\item Dibuje la gráfica del desplazamiento en función del tiempo.
\item Justifique la forma de la gráfica.
\end{enumerate}
